
\chapter{Implementation And Results}
\section{Software And Hardware Requirements}

The implementation of MetaGPT relies on a modern full-stack environment with the following software and hardware requirements:
\begin{itemize}
    \item \textbf{Hardware}: A development machine with a multi-core processor, at least 4~GB of RAM, and sufficient disk space to store generated projects. For typical experiments, a standard laptop or desktop configuration is sufficient.
    \item \textbf{Operating System}: Any operating system capable of running Python 3.11 and Node.js 18 or above (e.g., Windows, Linux, or macOS).
    \item \textbf{Backend Software}: Python 3.11+, FastAPI, LangChain, LangGraph, and supporting libraries managed using the \texttt{uv} package manager or \texttt{pip}. A valid Google API key is required to access the Gemini model.
    \item \textbf{Frontend Software}: Node.js 18+, Next.js 14, TypeScript, Tailwind CSS, and associated tooling such as ESLint and TypeScript for linting and type checking.
    \item \textbf{Tools}: Git for version control, a terminal or shell for running backend and frontend servers, and an editor or IDE for browsing code and logs.
\end{itemize}

\section{Implementation Details}

\subsection*{Module 1: Backend API Layer Implementation}

The backend is organized under \texttt{backend/\allowbreak{}app} and exposes its functionality through a FastAPI application defined in \texttt{main.py}. In this file, the FastAPI instance is created, CORS and other middleware are configured, and the top-level router from \texttt{api/\allowbreak{}router.py} is included, as shown in Listing~\ref{lst:fastapi_setup}. The router groups all versioned endpoints under the \texttt{/\allowbreak{}api/\allowbreak{}v1} prefix and delegates to modules in \texttt{api/\allowbreak{}endpoints}.

\begin{center}
\begin{minipage}{\linewidth}
\begin{lstlisting}[language=Python, caption={FastAPI application configuration and middleware setup}, label={lst:fastapi_setup}]
def create_app() -> FastAPI:
    settings = get_settings()
    app = FastAPI(
        title=settings.app_name,
        version=settings.app_version,
        lifespan=lifespan,
    )
    # Configure CORS
    app.add_middleware(
        CORSMiddleware,
        allow_origins=settings.cors_origins,
        allow_credentials=True,
        allow_methods=["*"],
        allow_headers=["*"],
    )
    # Include API routes
    app.include_router(api_router, prefix=settings.api_prefix)
    return app
\end{lstlisting}
\end{minipage}
\end{center}

Listing~\ref{lst:backend_create_project} shows a representative endpoint for creating a new project. It validates user credits, calls the pipeline service, and handles persistence.

\begin{center}
\begin{minipage}{\linewidth}
\begin{lstlisting}[language=Python, caption={FastAPI endpoint for project creation}, label={lst:backend_create_project}]
@router.post("", response_model=Project, status_code=status.HTTP_201_CREATED)
async def create_project(
    request: ProjectCreate,
    user: User = Depends(get_current_user),
) -> Project:
    if not user.has_credits():
        raise HTTPException(
            status_code=status.HTTP_403_FORBIDDEN,
            detail="You've used all your free credits. Upgrade to premium.",
        )
    service = PipelineService()
    project = await service.create_project(request.prompt, user_id=str(user.id))
    user.credits_used += 1
    await user.save()
    return project
\end{lstlisting}
\end{minipage}
\end{center}

The implementation of individual endpoints is split into focused modules:
\begin{itemize}
    \item \texttt{api/\allowbreak{}endpoints/\allowbreak{}projects.py} implements handlers for creating new projects, listing existing projects, and retrieving project metadata. It uses request/response models from \texttt{schemas/\allowbreak{}projects.py}.
    \item \texttt{api/\allowbreak{}endpoints/\allowbreak{}pipeline.py} contains endpoints to start a pipeline run (both standard and streaming) and to query pipeline status or artifacts, delegating to functions in \texttt{services/\allowbreak{}pipeline\_\allowbreak{}service.py}.
    \item \texttt{api/\allowbreak{}endpoints/\allowbreak{}files.py} provides a tree view of project files and operations to read or update file contents via the storage layer (\texttt{storage/\allowbreak{}file\_\allowbreak{}store.py}).
    \item \texttt{api/\allowbreak{}endpoints/\allowbreak{}chat.py} implements chat-based iteration endpoints that forward messages to \texttt{services/\allowbreak{}chat\_\allowbreak{}service.py}.
    \item \texttt{api/\allowbreak{}endpoints/\allowbreak{}rag.py} and \texttt{api/\allowbreak{}endpoints/\allowbreak{}sandbox.py} expose retrieval and sandbox-related operations.
\end{itemize}

Cross-cutting concerns such as authentication and authorization are implemented in \texttt{auth.py} and reused across endpoints using FastAPI dependencies. Configuration values required by the API layer are loaded from \texttt{config.py}, which reads from environment variables and provides typed settings via Pydantic.

\textbf{E2B Sandbox preview endpoints:} The sandbox endpoint module (\texttt{api/\allowbreak{}endpoints/\allowbreak{}sandbox.py}) is designed for long-running preview builds. It implements a non-blocking create flow that returns \textbf{202 Accepted} immediately and starts a background task that writes generated files, installs dependencies, and starts dev servers inside an isolated E2B sandbox. Users (via the frontend) poll the status endpoint until the sandbox becomes \texttt{ready}, at which point a public \texttt{preview\_\allowbreak{}url} is returned for embedding in an iframe.

\subsection*{Module 2: Agent Orchestration and Graph Implementation}

Agent orchestration is implemented in the \texttt{graph} and \texttt{agents} packages. The central pipeline state is defined in \texttt{graph/\allowbreak{}state.py}, where a Pydantic model captures the user prompt, derived requirements, architectural plan, generated files, and QA feedback. This state object is passed between nodes in a LangGraph graph.

The graph itself is constructed in \texttt{graph/\allowbreak{}orchestrator.py}. Here, nodes corresponding to the Manager, Architect, Engineer, and QA agents are registered, and edges are defined to specify the order of execution, as demonstrated in Listing~\ref{lst:orchestrator_graph}. Each node invokes a concrete agent class from the \texttt{agents} package:

\begin{center}
\begin{minipage}{\linewidth}
\begin{lstlisting}[language=Python, caption={LangGraph orchestrator pipeline configuration}, label={lst:orchestrator_graph}]
def _build_graph(self) -> StateGraph:
    workflow = StateGraph(PipelineState)
    # Add nodes for each agent
    workflow.add_node("manager", self._run_manager)
    workflow.add_node("architect", self._run_architect)
    workflow.add_node("engineer", self._run_engineer)
    workflow.add_node("qa", self._run_qa)

    # Define deterministic linear flow
    workflow.set_entry_point("manager")
    workflow.add_edge("manager", "architect")
    workflow.add_edge("architect", "engineer")
    workflow.add_edge("engineer", "qa")
    workflow.add_edge("qa", END)
    
    return workflow.compile()
\end{lstlisting}
\end{minipage}
\end{center}
\begin{itemize}
    \item \texttt{agents/\allowbreak{}manager.py} receives the raw prompt and produces structured requirements.
    \item \texttt{agents/\allowbreak{}architect.py} designs the system architecture and file layout based on those requirements.
    \item \texttt{agents/\allowbreak{}engineer.py} generates implementation code for the planned files. Listing~\ref{lst:engineer_prompt} illustrates how the Engineer Agent builds its LLM prompt using outputs from previous stages and SOP constraints.
    \item \texttt{agents/\allowbreak{}qa.py} inspects artifacts and produces test cases and quality feedback.
\end{itemize}

\begin{center}
\begin{minipage}{\linewidth}
\begin{lstlisting}[language=Python, caption={Engineer Agent building its prompt with context from earlier pipeline stages}, label={lst:engineer_prompt}]
def _build_prompt(self, **inputs: Any) -> str:
    architect_output: ArchitectOutput = inputs.get("architect_output")
    manager_output: ManagerOutput = inputs.get("manager_output")

    # Format SOP sections
    constraints_text = self._format_sop_section(self.sop.constraints)
    quality_text = self._format_sop_section(self.sop.quality_checklist)

    # Format architecture details
    file_structure = self._format_file_structure(architect_output)
    components = self._format_components(architect_output)
    requirements = self._format_requirements(manager_output)

    return f"""# Engineer Agent - Code Generation
## Your Role
{self.sop.role}

## Constraints
{constraints_text}

## Project: {manager_output.project_name}
{manager_output.project_description}

## File Structure:
{file_structure}

## Your Task
Implement ALL files listed in the file structure with production-ready code.
"""
\end{lstlisting}
\end{minipage}
\end{center}

All agents inherit from a base class in \texttt{agents/\allowbreak{}base.py}, which standardizes how LLM calls are made, how inputs/outputs are validated using \texttt{schemas/\allowbreak{}agents.py}, and how logging is performed. The orchestrator integrates with LangChain’s LLM abstractions provided by \texttt{llm/\allowbreak{}gemini.py}, enabling both free-form and structured-output interactions. The final orchestration API is exposed to the rest of the application through \texttt{services/\allowbreak{}pipeline\_\allowbreak{}service.py}, which provides functions to start runs, stream intermediate events, and fetch stored artifacts.

\subsection*{Module 3: Storage and Project Management Implementation}

Storage concerns are encapsulated in the \texttt{models} and \texttt{storage} packages. Domain models in \texttt{models/\allowbreak{}project.py} and \texttt{models/\allowbreak{}user.py} describe how projects and users are represented in memory, including fields such as identifiers, names, prompts, and timestamps. These models are mirrored by Pydantic schemas in \texttt{schemas/\allowbreak{}projects.py} to ensure consistency between API payloads and internal representations.

Project-level persistence is handled by \texttt{storage/\allowbreak{}project\_\allowbreak{}store.py}. This module is responsible for creating project directories under the configured \texttt{PROJECTS\_\allowbreak{}DIR}, writing metadata files if needed, and enumerating existing projects. File-level operations are implemented in \texttt{storage/\allowbreak{}file\_\allowbreak{}store.py}, which provides functions to:
\begin{itemize}
    \item Walk a project directory and build a hierarchical tree structure consumed by the frontend \texttt{FileExplorer}.
    \item Read file contents from disk and return them as text for display or further processing.
    \item Apply write operations when agents or users modify files, ensuring that paths are validated and confined to the project root.
\end{itemize}

These storage modules are consumed by both the API endpoints (for serving file trees and contents) and the agent pipeline (for writing generated code to disk). By centralizing filesystem access in a small number of modules, the system can later be extended to support alternative backends such as databases or cloud storage with minimal changes elsewhere.

\textbf{MongoDB persistence:} In our framework, project state is stored in MongoDB. The database connection is initialized in \texttt{db.py} using Motor and Beanie, and document models such as \texttt{models/\allowbreak{}project.py} and \texttt{models/\allowbreak{}user.py} are registered with Beanie. The \texttt{ProjectStore} (\texttt{storage/\allowbreak{}project\_\allowbreak{}store.py}) provides CRUD operations to create projects, update pipeline status, and persist agent outputs (Manager/Architect/Engineer/QA). Listing~\ref{lst:mongo_create} shows how a new project document is initialized and persisted. This database-backed state allows the application to restore generated files to disk when ephemeral storage is lost (e.g., after container restart), ensuring that the file explorer and subsequent operations remain consistent.

\begin{center}
\begin{minipage}{\linewidth}
\begin{lstlisting}[language=Python, caption={Creating and persisting a project document using Beanie ODM}, label={lst:mongo_create}]
async def create(self, project_id: str, prompt: str, user_id: str = "") -> Project:
    doc = ProjectDocument(
        project_id=project_id,
        user_id=user_id,
        prompt=prompt,
        state=ProjectState(
            pipeline_status=PipelineStatus(
                stage=PipelineStage.PENDING,
                progress=0,
            )
        ),
        created_at=datetime.utcnow(),
        updated_at=datetime.utcnow(),
    )
    await doc.insert()
    return self._to_project(doc)
\end{lstlisting}
\end{minipage}
\end{center}

\subsection*{Module 4: Frontend User Interface Implementation}

The frontend implementation lives in the \texttt{frontend/\allowbreak{}src} directory and is based on the Next.js App Router. Routing and layout are defined in \texttt{src/\allowbreak{}app/\allowbreak{}layout.tsx} and the various page components:
\begin{itemize}
    \item \texttt{src/\allowbreak{}app/\allowbreak{}page.tsx} implements the landing page where users can provide prompts and see high-level information about the system. A snippet handling the submission of a new project prompt is shown in Listing~\ref{lst:frontend_submit}.
    \item \texttt{src/\allowbreak{}app/\allowbreak{}projects/\allowbreak{}page.tsx} fetches the list of projects from the backend and renders them in a tabular or card-based layout.
    \item \texttt{src/\allowbreak{}app/\allowbreak{}project/\allowbreak{}[id]/\allowbreak{}page.tsx} is the main workspace, composed of panels for the file explorer, code viewer, agent outputs, execution timeline, and chat.
\end{itemize}

\begin{center}
\begin{minipage}{\linewidth}
\begin{lstlisting}[language=JavaScript, caption={Client-side form submission handler in Next.js}, label={lst:frontend_submit}]
const handleSubmit = async (e: React.FormEvent) => {
    e.preventDefault();
    if (!prompt.trim() || isLoading) return;
    if (!token) {
        router.push("/signin");
        return;
    }
    setIsLoading(true);
    try {
        const { createProject } = await import("@/lib/api");
        const project = await createProject(prompt.trim());
        router.push(`/project/${project.id}`);
    } catch (error: unknown) {
        setSubmitError("Failed to create project");
        setIsLoading(false);
    }
};
\end{lstlisting}
\end{minipage}
\end{center}

Key interactive components are implemented in \texttt{src/\allowbreak{}components}:
\begin{itemize}
    \item \texttt{FileExplorer.tsx} renders the tree returned by the files API and triggers file selection events. When only a flat list of generated files is available (e.g., streaming in real-time), it dynamically reconstructs the hierarchical folder structure as demonstrated in Listing~\ref{lst:frontend_file_tree}.
    \item \texttt{CodeViewer.tsx} uses a code editor or viewer (e.g., Monaco) to display and possibly edit source files.
    \item \texttt{AgentOutputs.tsx} subscribes to pipeline events and presents agent reasoning and outputs in a structured way.
    \item \texttt{ExecutionTimeline.tsx} visualizes the order and status of agent executions for a given run.
    \item \texttt{ChatPanel.tsx} provides a text area and message list for iterative refinement of a project as shown handling submissions in Listing~\ref{lst:chat_submit}.
\end{itemize}

\begin{center}
\begin{minipage}{\linewidth}
\begin{lstlisting}[language=JavaScript, caption={Dynamically building a file tree structure from a flat list of file paths}, label={lst:frontend_file_tree}]
function buildTreeFromFiles(files: { file_path: string }[]): FileTreeNode | null {
  if (files.length === 0) return null;

  const root: FileTreeNode = { name: "files", path: "/", type: "directory", children: [], };

  for (const file of files) {
    const parts = file.file_path.split("/").filter(Boolean);
    let current = root;

    for (let i = 0; i < parts.length; i++) {
      const part = parts[i];
      const isFile = i === parts.length - 1;
      const path = parts.slice(0, i + 1).join("/");
      let child = current.children?.find((c) => c.name === part);

      if (!child) {
        child = { name: part, path, type: isFile ? "file" : "directory", children: isFile ? [] : [], };
        current.children = current.children || [];
        current.children.push(child);
      }
      current = child;
    }
  }
  return root;
}
\end{lstlisting}
\end{minipage}
\end{center}

The module \texttt{src/\allowbreak{}lib/\allowbreak{}api.ts} centralizes HTTP calls to the backend, while \texttt{src/\allowbreak{}lib/\allowbreak{}store.ts} and related files implement client-side state management (e.g., current project, selected file, pipeline status). Tailwind CSS classes and global styles in \texttt{src/\allowbreak{}app/\allowbreak{}globals.css} are used to implement the dark theme and responsive layout.


\begin{center}
\begin{minipage}{\linewidth}
\begin{lstlisting}[language=JavaScript, caption={Chat handler supporting Server-Sent Events (SSE) and background indexing}, label={lst:chat_submit}]
const handleSubmit = async (e: React.FormEvent) => {
    e.preventDefault();
    if (!message.trim() || isLoading) return;

    const userMessage = message.trim();
    setMessage("");
    setIsLoading(true);

    try {
        const response = await sendChatMessage(
            projectId,
            userMessage,
            autoMode ? null : selectedModel,
        );

        addChatMessage({
            ...response.message,
            files_referenced: response.files_referenced,
            indexing_status: response.indexing_status,
        });

        // If indexing is pending, trigger it live
        if (response.indexing_status === "pending") {
            indexProjectFiles(projectId, response.files_modified);
        }
    } finally {
        setIsLoading(false);
    }
};
\end{lstlisting}
\end{minipage}
\end{center}


\subsection*{Module 5: LLM, SOP, and Retrieval Implementation}

The LLM, SOP, and retrieval logic is implemented in the \texttt{llm}, \texttt{sop}, and \texttt{rag} packages. The file \texttt{llm/\allowbreak{}gemini.py} exposes functions such as \texttt{get\_\allowbreak{}llm()} and \texttt{get\_\allowbreak{}llm\_\allowbreak{}with\_\allowbreak{}structured\_\allowbreak{}output(schema)} (see Listing~\ref{lst:gemini_structured}). These functions read configuration from \texttt{config.py} (e.g., model name, temperature, API key) and return LangChain \texttt{ChatModel} instances configured for MetaGPT. This ensures that all agents share a single, centrally managed LLM configuration.

\begin{center}
\begin{minipage}{\linewidth}
\begin{lstlisting}[language=Python, caption={Configuring Gemini for structured JSON output via LangChain}, label={lst:gemini_structured}]
def get_llm_with_structured_output(
    output_schema: type[T],
    temperature: float | None = None,
    max_tokens: int | None = None,
    model: str | None = None,
) -> BaseChatModel:
    base_llm = get_llm(temperature=temperature, max_tokens=max_tokens, model=model)
    return base_llm.with_structured_output(output_schema)
\end{lstlisting}
\end{minipage}
\end{center}

Agent SOPs are defined in \texttt{sop/\allowbreak{}definitions.py}. Each SOP includes fields such as role description, objectives, required inputs, expected outputs, constraints, and a quality checklist. The concrete agent classes in \texttt{agents/\allowbreak{}*.py} load these SOPs and use them to construct prompts or system messages when calling the LLM, ensuring that behavior is consistent and traceable across runs.

Retrieval-augmented generation is supported by the \texttt{rag} package:
\begin{itemize}
    \item \texttt{rag/\allowbreak{}embeddings.py} configures embedding models used to convert text into vector representations.
    \item \texttt{rag/\allowbreak{}vector\_\allowbreak{}store.py} manages storage and querying of embeddings in a vector database or in-memory index.
    \item \texttt{rag/\allowbreak{}indexer.py} defines workflows for indexing new documents, such as project files or external resources.
    \item \texttt{rag/\allowbreak{}retriever.py} implements retrieval logic that selects the most relevant chunks of text to be added as context to LLM prompts. Listing~\ref{lst:retriever_search} highlights the core similarity search function.
\end{itemize}

\begin{center}
\begin{minipage}{\linewidth}
\begin{lstlisting}[language=Python, caption={Retrieving relevant codebase chunks from the vector store}, label={lst:retriever_search}]
async def retrieve(self, project_id: str, query: str, k: int | None = None) -> list[RetrievedChunk]:
    if not self.settings.rag_enabled:
        return []

    k = k or self.settings.rag_top_k
    try:
        vector_store = get_vector_store(project_id)
        # Similarity search with scores
        results = vector_store.similarity_search_with_relevance_scores(query, k=k)

        chunks = []
        for doc, score in results:
            chunks.append(
                RetrievedChunk(
                    content=doc.page_content,
                    file_path=doc.metadata.get("file_path", "unknown"),
                    language=doc.metadata.get("language", "text"),
                    relevance_score=score,
                )
            )
        return chunks
    except Exception as e:
        logger.error(f"Retrieval error for project {project_id}: {e}")
        return []
\end{lstlisting}
\end{minipage}
\end{center}

By combining these pieces, the system can ground agent reasoning in both the current project state and external knowledge sources, improving the relevance and reliability of generated code and explanations.

\section{Testing}
\subsection{Unit Testing}

Table~\ref{tab:unit_tests} summarizes representative unit test cases for the main modules of MetaGPT. Each test case specifies the module under test, the scenario being validated, sample input, expected result, and the final status of execution.

\renewcommand{\arraystretch}{1.3}
\begingroup
\setlength{\LTleft}{-1.5cm}
\begin{longtable}{|p{3.0cm}|p{3.8cm}|p{3.8cm}|p{4.6cm}|>{\centering\arraybackslash}p{1.8cm}|}
\caption{Unit Test Cases for MetaGPT Modules}
\label{tab:unit_tests}
\\
\hline
\textbf{Module} & \textbf{Test Scenario} & \textbf{Input} & \textbf{Expected Result} & \textbf{Status} \\
\hline
\endfirsthead
\hline
\textbf{Module} & \textbf{Test Scenario} & \textbf{Input} & \textbf{Expected Result} & \textbf{Status} \\
\hline
\endhead
\hline
\endfoot
\hline
\endlastfoot
Backend API Layer & Load configuration with valid environment variables & Populated \texttt{.env} file (e.g., \texttt{GOOGLE\_\allowbreak{}API\_\allowbreak{}KEY}, \texttt{PROJECTS\_\allowbreak{}DIR}) & \texttt{Config} object created successfully with all fields set and no validation errors & Success \\
\hline
Backend API Layer & Handle missing optional configuration values & \texttt{.env} file without optional fields (e.g., no \texttt{LLM\_\allowbreak{}TEMPERATURE}) & Default values applied for missing options; application starts without raising exceptions & Success \\
\hline
Backend API Layer & Validate project creation endpoint & HTTP \texttt{POST} \texttt{/\allowbreak{}api/\allowbreak{}v1/\allowbreak{}projects} with a valid JSON body containing a prompt & Response status \texttt{201}; response body contains non-empty project ID and stored prompt & Success \\
\hline
Backend API Layer & Reject invalid project request payload & HTTP \texttt{POST} \texttt{/\allowbreak{}api/\allowbreak{}v1/\allowbreak{}projects} with missing required fields & Response status \texttt{422} (validation error); descriptive error message returned & Success \\
\hline
Backend API Layer & Fetch project list when no projects exist & HTTP \texttt{GET} \texttt{/\allowbreak{}api/\allowbreak{}v1/\allowbreak{}projects} on a clean workspace & Response status \texttt{200}; response body is an empty list & Success \\
\hline
Agent Orchestration and Graph & Initialize empty pipeline state & Call state constructor in \texttt{graph/\allowbreak{}state.py} with a sample prompt & State object created with prompt set and other fields initialized to default/empty values & Success \\
\hline
Agent Orchestration and Graph & Manager agent produces structured requirements & Invoke Manager agent with a simple prompt (mocked LLM) & Output conforms to \texttt{ManagerOutput} schema and contains non-empty requirement fields & Success \\
\hline
Agent Orchestration and Graph & Architect agent generates file plan & Provide Manager output to Architect agent (mocked LLM) & Output includes a file list with valid relative paths and descriptions & Success \\
\hline
Agent Orchestration and Graph & Graph orchestrator runs full pipeline with mocked agents & Execute orchestrator with all agents mocked to return fixed outputs & Final state contains populated requirements, architecture, code artifacts, and QA data without errors & Success \\
\hline
Agent Orchestration and Graph & Error handling in a single node & Force an exception inside one agent node & Orchestrator catches error, records it in state/logs, and terminates or retries according to policy & Success \\
\hline
Storage and Project Management & Create new project directory & Call function in \texttt{project\_\allowbreak{}store} with a new project ID & Corresponding folder created under \texttt{PROJECTS\_\allowbreak{}DIR}; path is returned and exists on disk & Success \\
\hline
Storage and Project Management & Build file tree for simple project & Populate a test project with a few files and call \texttt{file\_\allowbreak{}store} tree builder & Returned tree structure matches directory layout and contains correct file names and types & Success \\
\hline
Storage and Project Management & Read file contents safely & Request contents of a known file via \texttt{file\_\allowbreak{}store} & Function returns exact file text; attempts to escape project root are rejected & Success \\
\hline
Storage and Project Management & Update file content & Call update function with a valid file path and new code content & File on disk is overwritten with new content; subsequent read returns updated text & Success \\
\hline
Storage and Project Management & Handle non-existent file read & Request a file that does not exist in the project tree & Function raises a controlled error or returns a standardized ``not found'' response & Success \\
\hline
Frontend User Interface & Render project list page with data & Mock API client to return a list of projects and render \texttt{projects/\allowbreak{}page.tsx} & Page displays the correct number of project entries with expected titles and IDs & Success \\
\hline
Frontend User Interface & Display file contents in code viewer & Provide a selected file and mocked content to \texttt{CodeViewer.tsx} & Component renders the content in the editor area without runtime errors & Success \\
\hline
Frontend User Interface & Update selected file state on tree click & Simulate click on a file node in \texttt{FileExplorer.tsx} & Global/store state updates to mark that file as selected; dependent components receive the update & Success \\
\hline
Frontend User Interface & Show execution timeline events & Supply a sequence of mocked pipeline events to \texttt{ExecutionTimeline.tsx} & Timeline renders entries in correct order with appropriate status indicators & Success \\
\hline
Frontend User Interface & Send chat message to backend & Trigger a send action in \texttt{ChatPanel.tsx} with mocked API client & API client is called with correct project ID and message payload; UI input is cleared & Success \\
\hline
LLM, SOP, and Retrieval & Construct base LLM instance & Call \texttt{get\_\allowbreak{}llm()} from \texttt{llm/\allowbreak{}gemini.py} with valid configuration & Function returns a non-\texttt{None} LangChain model instance using configured parameters & Success \\
\hline
LLM, SOP, and Retrieval & Create structured-output LLM & Call \texttt{get\_\allowbreak{}llm\_\allowbreak{}with\_\allowbreak{}}\linebreak\texttt{structured\_\allowbreak{}output()} with a schema & Returned model enforces the specified schema and raises an error for invalid outputs in tests & Success \\
\hline
LLM, SOP, and Retrieval & Load SOP definitions & Import and access SOPs from \texttt{sop/\allowbreak{}definitions.py} & Each defined agent SOP contains non-empty role, objective, inputs, outputs, and checklist fields & Success \\
\hline
LLM, SOP, and Retrieval & Generate embeddings for sample text & Call embedding function in \texttt{rag/\allowbreak{}embeddings.py} with a short string & Function returns a numeric vector of expected dimension without errors & Success \\
\hline
LLM, SOP, and Retrieval & Retrieve nearest neighbor from vector store & Index two or more documents and query via \texttt{rag/\allowbreak{}retriever.py} & Retriever returns the most similar document to the query text according to test setup & Success \\
\hline
\end{longtable}
\endgroup

\subsection{Integration Testing}

Integration testing verifies that the various subsystems of MetaGPT---FastAPI endpoints, agent pipeline, storage, and RAG utilities---work together correctly when exercised through realistic workflows. Tables~\ref{tab:integration_api}--\ref{tab:integration_rag} summarize integration test cases for each area, using the same template as the unit testing table.

\subsubsection*{API-level Tests}
Using \texttt{pytest} and FastAPI's \texttt{TestClient}, integration tests call endpoints defined in \texttt{backend/\allowbreak{}app/\allowbreak{}api/\allowbreak{}endpoints}. Table~\ref{tab:integration_api} lists representative scenarios.

\renewcommand{\arraystretch}{1.2}
\begin{longtable}{|p{2.8cm}|p{3.4cm}|p{3.4cm}|p{4.2cm}|>{\centering\arraybackslash}p{1.6cm}|}
\caption{Integration Test Cases: API-level Tests}
\label{tab:integration_api}
\\
\hline
\textbf{Module} & \textbf{Test Scenario} & \textbf{Input} & \textbf{Expected Result} & \textbf{Status} \\
\hline
\endfirsthead
\hline
\textbf{Module} & \textbf{Test Scenario} & \textbf{Input} & \textbf{Expected Result} & \textbf{Status} \\
\hline
\endhead
\hline
\endfoot
\hline
\endlastfoot
API (Projects) & Create project and verify persistence & POST \texttt{/\allowbreak{}api/\allowbreak{}v1/\allowbreak{}projects} with \texttt{\{"prompt": "Todo app"\}} & Status 201; response has project \texttt{id}; project dir exists under \texttt{PROJECTS\_\allowbreak{}DIR} & Success \\
\hline
API (Projects) & List projects after creating one & GET \texttt{/\allowbreak{}api/\allowbreak{}v1/\allowbreak{}projects} after creating a project & Status 200; list contains at least one project with matching \texttt{id} and prompt & Success \\
\hline
API (Pipeline) & Start pipeline run and get status & POST \texttt{/\allowbreak{}api/\allowbreak{}v1/\allowbreak{}pipeline/\allowbreak{}run} with valid project \texttt{id}; then GET \texttt{/\allowbreak{}api/\allowbreak{}v1/\allowbreak{}pipeline/\allowbreak{}\{id\}/status} & Run starts; status endpoint returns run state (e.g., running/completed) and artifact keys & Success \\
\hline
API (Pipeline) & Stream pipeline events & POST \texttt{/\allowbreak{}api/\allowbreak{}v1/\allowbreak{}pipeline/\allowbreak{}stream} with prompt; consume SSE stream & Stream yields events (e.g., manager\_done, architect\_done); connection stays open until completion & Success \\
\hline
API (Files) & Get file tree for generated project & GET \texttt{/\allowbreak{}api/\allowbreak{}v1/\allowbreak{}files/\allowbreak{}\{id\}/tree} for a project with generated files & Status 200; JSON tree matches on-disk structure; nodes have \texttt{name} and \texttt{type} (file/dir) & Success \\
\hline
API (Files) & Read file content via API & GET \texttt{/\allowbreak{}api/\allowbreak{}v1/\allowbreak{}files/\allowbreak{}\{id\}/content/\{path\}} for a known file & Status 200; body contains exact file content; content-type appropriate for text & Success \\
\hline
API (Chat) & Send chat message and fetch history & POST \texttt{/\allowbreak{}api/\allowbreak{}v1/\allowbreak{}chat/\allowbreak{}\{id\}} with message; GET \texttt{/\allowbreak{}api/\allowbreak{}v1/\allowbreak{}chat/\allowbreak{}\{id\}/history} & Message stored; history returns list including the new message with correct \texttt{role} and \texttt{content} & Success \\
\hline
\end{longtable}

\subsubsection*{Pipeline Orchestration}
Tests exercise the full agent pipeline (Manager $\rightarrow$ Architect $\rightarrow$ Engineer $\rightarrow$ QA) via the LangGraph orchestrator. Table~\ref{tab:integration_pipeline} lists key scenarios.

\renewcommand{\arraystretch}{1.2}
\begin{longtable}{|p{2.8cm}|p{3.4cm}|p{3.4cm}|p{4.2cm}|>{\centering\arraybackslash}p{1.6cm}|}
\caption{Integration Test Cases: Pipeline Orchestration}
\label{tab:integration_pipeline}
\\
\hline
\textbf{Module} & \textbf{Test Scenario} & \textbf{Input} & \textbf{Expected Result} & \textbf{Status} \\
\hline
\endfirsthead
\hline
\textbf{Module} & \textbf{Test Scenario} & \textbf{Input} & \textbf{Expected Result} & \textbf{Status} \\
\hline
\endhead
\hline
\endfoot
\hline
\endlastfoot
Pipeline (Graph) & Full pipeline run with mocked LLM & Single prompt; all agents use mocked LLM returning valid structured outputs & Final state has requirements, architecture, file list, generated code, and QA output; no unhandled exceptions & Success \\
\hline
Pipeline (Graph) & State passed correctly between nodes & Inspect state after each node (Manager, Architect, Engineer, QA) & Each node receives correct input from previous node; state fields updated in order (requirements $\rightarrow$ design $\rightarrow$ code $\rightarrow$ qa) & Success \\
\hline
Pipeline (Services) & Pipeline service starts run and returns run ID & Call \texttt{pipeline\_\allowbreak{}service} start with project ID and prompt & Returns run identifier; orchestrator invoked; state persisted or streamed as configured & Success \\
\hline
Pipeline (Services) & Retrieve artifacts after run & GET artifacts for a completed pipeline run ID & Response includes manager output, architect output, engineer outputs (files), QA output in expected structure & Success \\
\hline
Pipeline (Graph) & Error in one agent does not corrupt state & Force one agent to raise; capture state/logs & Error caught; state or logs indicate which node failed; no partial writes leave inconsistent state & Success \\
\hline
\end{longtable}

\subsubsection*{Storage + API Integration}
End-to-end tests combine storage and API: create project, run pipeline (or seed files), then access via \texttt{/\allowbreak{}api/\allowbreak{}v1/\allowbreak{}files}. Table~\ref{tab:integration_storage} summarizes test cases.

\renewcommand{\arraystretch}{1.2}
\begin{longtable}{|p{2.8cm}|p{3.4cm}|p{3.4cm}|p{4.2cm}|>{\centering\arraybackslash}p{1.6cm}|}
\caption{Integration Test Cases: Storage + API Integration}
\label{tab:integration_storage}
\\
\hline
\textbf{Module} & \textbf{Test Scenario} & \textbf{Input} & \textbf{Expected Result} & \textbf{Status} \\
\hline
\endfirsthead
\hline
\textbf{Module} & \textbf{Test Scenario} & \textbf{Input} & \textbf{Expected Result} & \textbf{Status} \\
\hline
\endhead
\hline
\endfoot
\hline
\endlastfoot
Storage + API & Create project then list via API & Create project via API; add a file to project dir (simulate pipeline); GET \texttt{/\allowbreak{}api/\allowbreak{}v1/\allowbreak{}files/\allowbreak{}\{id\}/tree} & Tree includes the new file; path and name match filesystem & Success \\
\hline
Storage + API & Read file written by pipeline & Run pipeline (or seed file); GET \texttt{/\allowbreak{}api/\allowbreak{}v1/\allowbreak{}files/\allowbreak{}\{id\}/content/\allowbreak{}\{path\}} & Content returned matches content written by pipeline or test & Success \\
\hline
Storage + API & Update file via API and read back & PUT \texttt{/\allowbreak{}api/\allowbreak{}v1/\allowbreak{}files/\allowbreak{}\{id\}/content/\allowbreak{}\{path\}} with new body; GET same path & File on disk updated; GET returns new content & Success \\
\hline
Storage + API & Invalid path returns error & GET \texttt{/\allowbreak{}api/\allowbreak{}v1/\allowbreak{}files/\allowbreak{}\{id\}/content/\allowbreak{}\{path\}/../etc/\allowbreak{}passwd} or path outside project & Status 400/403/404; no content from outside project root returned & Success \\
\hline
Storage + API & Empty project returns empty tree & GET \texttt{/\allowbreak{}api/\allowbreak{}v1/\allowbreak{}files/\allowbreak{}\{id\}/tree} for newly created project with no files & Status 200; tree is empty or has only root node & Success \\
\hline
\end{longtable}

\subsubsection*{RAG Workflows}
For setups where RAG is enabled, integration tests index documents and call retrieval endpoints. Table~\ref{tab:integration_rag} lists representative cases.

\renewcommand{\arraystretch}{1.2}
\begin{longtable}{|p{2.8cm}|p{3.4cm}|p{3.4cm}|p{4.2cm}|>{\centering\arraybackslash}p{1.6cm}|}
\caption{Integration Test Cases: RAG Workflows}
\label{tab:integration_rag}
\\
\hline
\textbf{Module} & \textbf{Test Scenario} & \textbf{Input} & \textbf{Expected Result} & \textbf{Status} \\
\hline
\endfirsthead
\hline
\textbf{Module} & \textbf{Test Scenario} & \textbf{Input} & \textbf{Expected Result} & \textbf{Status} \\
\hline
\endhead
\hline
\endfoot
\hline
\endlastfoot
RAG (Indexer) & Index small document set & Call indexer with a few code snippets or file paths & Index completes; vector store contains expected number of entries; no errors & Success \\
\hline
RAG (Retriever) & Query returns relevant chunks & Index 2--3 documents; query with text similar to one document & Retriever returns chunks; top result matches or is highly similar to the queried document & Success \\
\hline
RAG + API & Retrieval endpoint returns results & GET/POST RAG endpoint with project ID and query string & Status 200; response includes retrieved snippets or IDs usable for agent context & Success \\
\hline
RAG (Embeddings) & Embeddings consistent for same text & Call embedding function twice with same string & Both calls return identical (or numerically equal) vector & Success \\
\hline
RAG (Vector store) & Add and query in same session & Add documents to vector store; run query & Query returns added documents when query text is relevant; no cross-session persistence errors & Success \\
\hline
\end{longtable}

\subsubsection*{Integration of Backend to Frontend}
Tests verify that the Next.js frontend correctly communicates with the FastAPI backend: API base URL, request/response handling, and real-time updates. Table~\ref{tab:integration_backend_frontend} lists representative scenarios.

\renewcommand{\arraystretch}{1.2}
\begin{longtable}{|p{2.8cm}|p{3.4cm}|p{3.4cm}|p{4.2cm}|>{\centering\arraybackslash}p{1.6cm}|}
\caption{Integration Test Cases: Backend to Frontend}
\label{tab:integration_backend_frontend}
\\
\hline
\textbf{Module} & \textbf{Test Scenario} & \textbf{Input} & \textbf{Expected Result} & \textbf{Status} \\
\hline
\endfirsthead
\hline
\textbf{Module} & \textbf{Test Scenario} & \textbf{Input} & \textbf{Expected Result} & \textbf{Status} \\
\hline
\endhead
\hline
\endfoot
\hline
\endlastfoot
Frontend + API & Projects list loads from backend & Open \texttt{/\allowbreak{}projects}; frontend calls GET \texttt{/\allowbreak{}api/\allowbreak{}v1/\allowbreak{}projects} & Page displays project list; data matches backend response; no CORS or network errors & Success \\
\hline
Frontend + API & Create project from UI and see in list & Submit prompt on landing/workspace; frontend POSTs to \texttt{/\allowbreak{}api/\allowbreak{}v1/\allowbreak{}projects} or pipeline & New project appears in list or workspace; project ID and prompt reflected in UI & Success \\
\hline
Frontend + API & File tree and content from backend & Select project; frontend GET \texttt{/\allowbreak{}api/\allowbreak{}v1/\allowbreak{}files/\allowbreak{}\{id\}/tree} and content endpoints & FileExplorer shows tree; CodeViewer shows file content; data matches backend storage & Success \\
\hline
Frontend + API & Pipeline run and streaming updates & Start pipeline from UI; frontend uses \texttt{/\allowbreak{}api/\allowbreak{}v1/\allowbreak{}pipeline/\allowbreak{}stream} or status poll & Execution timeline updates; agent outputs appear; run completes or shows progress without timeout & Success \\
\hline
Frontend + API & Chat message sent and history loaded & Send message in ChatPanel; frontend POST \texttt{/\allowbreak{}api/\allowbreak{}v1/\allowbreak{}chat/\allowbreak{}\{id\}}; fetch history & Message appears in UI; history GET returns same message; no duplicate or lost messages & Success \\
\hline
Frontend + API & API base URL and env configuration & Run frontend with \texttt{NEXT\_\allowbreak{}PUBLIC\_\allowbreak{}API\_\allowbreak{}URL} pointing to backend & All API calls target correct host/port; no mixed content or wrong-origin errors & Success \\
\hline
\end{longtable}

These integration tests are run alongside unit tests with:
\begin{verbatim}
cd backend
uv run pytest
\end{verbatim}
and are often marked or grouped to distinguish slower, end-to-end scenarios from fast unit tests. For the frontend, integration is validated primarily by running the application against a local backend, while automated checks such as \texttt{npm run lint} and \texttt{npm run type-check} ensure consistency between API usage and TypeScript types.

\subsection{System Testing}

System testing validates the MetaGPT project as a whole from an end-user perspective: full workflow from prompt entry to generated project, UI responsiveness, and correctness of end-to-end behaviour. Table~\ref{tab:system_tests} summarizes system-level test cases.

\renewcommand{\arraystretch}{1.2}
\begin{longtable}{|p{2.8cm}|p{3.4cm}|p{3.4cm}|p{4.2cm}|>{\centering\arraybackslash}p{1.6cm}|}
\caption{System Test Cases: Project as a Whole}
\label{tab:system_tests}
\\
\hline
\textbf{Module} & \textbf{Test Scenario} & \textbf{Input} & \textbf{Expected Result} & \textbf{Status} \\
\hline
\endfirsthead
\hline
\textbf{Module} & \textbf{Test Scenario} & \textbf{Input} & \textbf{Expected Result} & \textbf{Status} \\
\hline
\endhead
\hline
\endfoot
\hline
\endlastfoot
End-to-end & Create project from prompt to codebase & User enters prompt (e.g., ``Build a todo app with React''); runs pipeline; waits for completion & Project created; pipeline runs Manager $\rightarrow$ Architect $\rightarrow$ Engineer $\rightarrow$ QA; generated files appear in file tree; no fatal errors & Success \\
\hline
End-to-end & View and browse generated project & User opens project workspace after pipeline completion & File tree shows generated structure; selecting a file displays content in code viewer; paths and content match backend & Success \\
\hline
End-to-end & Iterate via chat and re-run pipeline & User sends chat message requesting a change; triggers follow-up run or refinement & Chat history updated; pipeline or agent responds; project state remains consistent; UI reflects changes & Success \\
\hline
End-to-end & Multiple projects in sequence & User creates Project A, then Project B; lists projects; opens each & Both projects listed; switching between projects shows correct file tree and content for each; no data mixed between projects & Success \\
\hline
End-to-end & Execution timeline reflects pipeline progress & User starts pipeline; observes timeline during run & Timeline shows agent steps (Manager, Architect, Engineer, QA) in order; status indicators update; completion state shown & Success \\
\hline
End-to-end & Landing and navigation & User opens app; navigates Landing $\rightarrow$ Projects $\rightarrow$ Project workspace & All pages load; navigation works; no blank screens or console errors; backend reachable from frontend & Success \\
\hline
End-to-end & Error handling when backend unavailable & Stop backend; perform action from frontend (e.g., load projects) & Frontend shows error or fallback; no uncaught exception; user can retry or see clear message & Success \\
\hline
End-to-end & Generated project structure is usable & Inspect generated project on disk (e.g., \texttt{PROJECTS\_\allowbreak{}DIR}) & Directory structure exists; key files (e.g., package.json, main components) present; structure aligns with Architect output & Success \\
\hline
\end{longtable}

\section{Experimental Results \& Discussion}

To evaluate MetaGPT, a series of experiments were performed using representative prompts such as ``Build a todo app with React and a REST API backend'' and ``Build a blog platform with Next.js''. For each prompt, the agent pipeline was executed end-to-end, and the resulting project structure, code quality, and test coverage were examined. The system successfully generated multi-file codebases with clear separation between frontend and backend components, along with auxiliary files such as configuration and routing modules.

Qualitative analysis showed that the SOP-driven pipeline led to more consistent architectures and clearer traceability from requirements to implementation compared to ad-hoc single-prompt generation. The execution timeline and agent output visualization also improved transparency, allowing users to understand and debug the reasoning process of each agent. Limitations observed during experimentation include sensitivity to ambiguous prompts and occasional need for manual adjustments to align generated code with specific framework versions or organizational coding standards. These observations inform the future enhancements discussed in later chapters.
